\documentclass[11pt]{article}
\usepackage{../../../tex/lecturestyle}

\title{Polynomials over F\_p}
\author{}
\date{}

\begin{document}
\thispagestyle{fancy}

\lecturetitleband{Number Theory}{Polynomials over $\mathbb{F}_p$}{5}

\section{Overview}
Working with polynomials whose coefficients live in $\mathbb{F}_p$ (the
integers mod a prime $p$, under the field structure this gives them) is a
foundational tool in competitive programming, even though a problem
rarely names ``polynomials over $\mathbb{F}_p$'' explicitly. Most of what
follows is the same algebra taught for real- or complex-coefficient
polynomials in secondary school; the point of this lecture is to confirm
it survives the switch to $\mathbb{F}_p$, and to flag the one place it
does not: everything here relies on $\mathbb{F}_p$ being a field (every
nonzero element has a multiplicative inverse), which fails for a
composite modulus. Assumes modular arithmetic and basic properties of
primes.

\section{Division, Remainder, and Factor Theorems}
\begin{definitionbox}
For $f, g \in \mathbb{F}_p[x]$ with $g \ne 0$, there exist unique
$q, r \in \mathbb{F}_p[x]$ with $f = gq + r$ and $\deg r < \deg g$. Call
$q$ the \emph{quotient} and $r$ the \emph{remainder}.
\end{definitionbox}
Existence and uniqueness follow the same long-division argument as over
$\mathbb{R}$: at each step, the leading term of the current remainder is
cancelled using the leading coefficient of $g$, which is invertible
because $\mathbb{F}_p$ is a field. That invertibility is doing real work
--- see \S5 for what breaks when the modulus is composite.

\begin{theorembox}[Remainder Theorem]
For $f \in \mathbb{F}_p[x]$ and $a \in \mathbb{F}_p$, the remainder of
$f(x)$ divided by $x - a$ equals $f(a)$.
\end{theorembox}
\begin{proof}
Write $f(x) = (x-a)q(x) + r(x)$ with $\deg r < 1$, so $r$ is a constant.
Substituting $x = a$ gives $r = f(a)$.
\end{proof}

\begin{theorembox}[Factor Theorem]
If $f(a) = 0$ for $f \in \mathbb{F}_p[x]$, $a \in \mathbb{F}_p$, then
$f(x) = (x-a)g(x)$ for some $g \in \mathbb{F}_p[x]$.
\end{theorembox}
This is immediate from the remainder theorem: the remainder is $f(a)=0$,
so $x-a$ divides $f$ exactly.

\begin{theorembox}[Factor theorem, iterated]
If $a_1,\ldots,a_n \in \mathbb{F}_p$ are distinct and all are roots of
$f \in \mathbb{F}_p[x]$, then
$f(x) = (x-a_1)(x-a_2)\cdots(x-a_n)\,g(x)$ for some
$g \in \mathbb{F}_p[x]$.
\end{theorembox}
\begin{proof}
Induction on $n$. The base case is the factor theorem above. For the
step, factor out $(x-a_1)\cdots(x-a_n)$ by the inductive hypothesis,
leaving $f(x) = (x-a_1)\cdots(x-a_n)g(x)$. Substituting $x = a_{n+1}$: the
product of the $n$ nonzero factors $(a_{n+1}-a_i)$ times $g(a_{n+1})$
must be $0$, and a product of nonzero field elements is nonzero, so
$g(a_{n+1}) = 0$. Apply the factor theorem to $g$.
\end{proof}

\section{How Many Roots Can a Polynomial Have?}
\begin{theorembox}
A nonzero degree-$n$ polynomial in $\mathbb{F}_p[x]$ has at most $n$
roots in $\mathbb{F}_p$.
\end{theorembox}
\begin{proof}
Suppose $f$ has $n$ or more roots $a_1,\ldots,a_n$; we show it has
exactly these. By the iterated factor theorem,
$f(x) = f_n(x-a_1)\cdots(x-a_n)$ where $f_n \ne 0$ is $f$'s leading
coefficient. A product of nonzero factors is nonzero, so $f(x) = 0$ only
at $x \in \{a_1,\ldots,a_n\}$.
\end{proof}

\begin{notebox}
It is worth pausing on exactly where primality was used, since it is easy
to take this bound for granted. \S5 below works through what happens to
every theorem in this section when $\mathbb{F}_p$ is replaced by
$\mathbb{Z}/n\mathbb{Z}$ for composite $n$ --- division stops being
well-defined at all, and the root-count bound fails outright (e.g.\ $2x$
has two roots, $0$ and $2$, in $\mathbb{Z}/4\mathbb{Z}$, despite being
degree $1$).
\end{notebox}

\section{The Factorization of \texorpdfstring{$x^p - x$}{x\^{}p - x}}
\begin{theorembox}
In $\mathbb{F}_p[x]$, $\displaystyle x^p - x = \prod_{a \in \mathbb{F}_p}
(x - a)$.
\end{theorembox}
\begin{proof}
By Fermat's little theorem, every $a \in \mathbb{F}_p$ satisfies
$a^p \equiv a$, i.e.\ every element of $\mathbb{F}_p$ is a root of
$x^p-x$. That is $p$ distinct roots of a degree-$p$ polynomial, so the
root-count bound forces the full factorization (both sides are monic of
degree $p$ with the same $p$ roots).
\end{proof}

\begin{theorembox}[Wilson's Theorem]
For prime $p$, $(p-1)! \equiv -1 \pmod p$.
\end{theorembox}
\begin{proof}
Compare the coefficient of $x^1$ on both sides of the factorization
above. On the right, expanding $\prod_a (x-a)$, the $x^1$ coefficient is
$(-1)^{p-1}\!\!\sum$ of all $\binom{p-1}{1}$-fold products of the
$p-1$ nonzero elements leaving out one factor of $x$ times the constant
terms of the rest --- concretely, it is $(-1)^{p-1}$ times the elementary
symmetric polynomial of degree $p-1$ in $0,1,\ldots,p-1$, which (since
one of the roots is $0$) reduces to $(-1)^{p-1}$ times the product of the
nonzero elements, $(p-1)!$. On the left, $x^p - x$ has $x^1$ coefficient
$-1$. For odd $p$, $(-1)^{p-1}=1$, giving $(p-1)! \equiv -1$; $p=2$ checks
directly ($1! = 1 \equiv -1 \pmod 2$).
\end{proof}

\section{What Requires Primality: A Worked Comparison}
\begin{examplebox}
Replacing $\mathbb{F}_p$ with $\mathbb{Z}/n\mathbb{Z}$ for composite $n$
throughout this lecture:
\begin{itemize}[leftmargin=1.5em]
  \item \textbf{Division} (\S2, Definition): fails in general. In
        $\mathbb{Z}/4\mathbb{Z}$, no $q,r$ satisfy $x^2 = (2x)q(x)+r(x)$
        with $\deg r < 1$.
  \item \textbf{Remainder theorem}: division by a general $g$ is not even
        defined, but division by $x-a$ specifically still works (its
        leading coefficient is $1$, always invertible), so
        $f(x)=(x-a)q(x)+f(a)$ still holds in this restricted sense.
  \item \textbf{Factor theorem}: still holds --- it only used the
        remainder theorem with $g=x-a$.
  \item \textbf{Iterated factor theorem}: fails. In
        $\mathbb{Z}/4\mathbb{Z}$, $f(x)=2x$ has roots $a_1=0,a_2=2$, but
        $2x \ne (x-0)(x-2)g(x)$ for any $g$ (the right side is degree
        $\ge 2$). The induction step broke because
        $g(a_{n+1})=0$ no longer follows from a product of nonzero
        \emph{ring} elements being zero --- $\mathbb{Z}/4\mathbb{Z}$ has
        zero divisors ($2\cdot 2 \equiv 0$).
  \item \textbf{Root-count bound}: fails, by the same example ($2x$, a
        degree-$1$ polynomial, has $2$ roots in
        $\mathbb{Z}/4\mathbb{Z}$); a cleaner example with more roots than
        degree is $x^2-1$ in $\mathbb{Z}/8\mathbb{Z}$, with roots
        $1,3,5,7$.
\end{itemize}
The reference code below checks every one of these claims directly.
\end{examplebox}

\section{Schwartz--Zippel Lemma}
This is the multivariate generalization of the root-count bound, and a
standard ingredient in randomized-algorithm correctness proofs (identity
testing, hashing, and Tutte-matrix-based matching algorithms among them).

\begin{theorembox}[Schwartz--Zippel]
Let $f \in \mathbb{F}_p[x_1,\ldots,x_n]$ be a nonzero polynomial of total
degree $d$. Then the number of $n$-tuples $(a_1,\ldots,a_n) \in
\mathbb{F}_p^n$ with $f(a_1,\ldots,a_n)=0$ is at most $d\cdot p^{n-1}$.
Equivalently, for a uniformly random tuple,
$\Pr[f(a_1,\ldots,a_n)=0] \le d/p$.
\end{theorembox}
\begin{proof}
Induction on $n$; the base case $n=1$ is the root-count bound above.
For the step, write $f(x_1,\ldots,x_n,t)$ (naming the $(n{+}1)$-th
variable $t$) as a polynomial in $t$ with coefficients
$g_i \in \mathbb{F}_p[x_1,\ldots,x_n]$:
\[
  f(x_1,\ldots,x_n,t) = \sum_{i=0}^{k} g_i(x_1,\ldots,x_n)\, t^i,
\]
where $k=\deg_t f$. Since $f \ne 0$, $g_k \ne 0$, and $\deg g_k \le d-k$.
For a uniformly random $(a_1,\ldots,a_n,a_{n+1})$: by the inductive
hypothesis, $\Pr[g_k(a_1,\ldots,a_n)=0] \le (d-k)/p$; and conditioned on
$g_k(a_1,\ldots,a_n)\ne 0$, the base case applied to
$f(a_1,\ldots,a_n,t)$ (a nonzero degree-$k$ polynomial in $t$ alone)
gives $\Pr[f(a_1,\ldots,a_n,a_{n+1})=0] \le k/p$. Summing,
$\Pr[f=0] \le (d-k)/p + k/p = d/p$.
\end{proof}

\section{Reference Implementations}
The examples below implement polynomial long division over
$\mathbb{F}_p$ from scratch and check it, via a large batch of random
trials, against direct evaluation (the remainder theorem), then verify
the factor theorem and the root-count bound by finding roots through
brute-force evaluation (feasible since a small prime is used) and
re-expanding the resulting product of linear factors to confirm it
reproduces the original polynomial. They also expand
$\prod_{a\in\mathbb{F}_p}(x-a)$ directly and check the result against
$x^p-x$ coefficient-by-coefficient, cross-check the resulting Wilson's
theorem identity against a direct factorial computation, and verify
Schwartz--Zippel by exhaustively evaluating random bivariate polynomials
over all of $\mathbb{F}_p \times \mathbb{F}_p$ and measuring how often
they vanish. All checks use a small prime ($p=17$) specifically so that
``exhaustive over $\mathbb{F}_p$'' checks stay cheap.

\subsection{Rust}
\lstinputlisting[language=Rust]{example.rs}

\subsection{Python}
\lstinputlisting[language=Python]{example.py}

\section{Worked Example}
\begin{examplebox}
With $p=17$: the remainder theorem, factor theorem, and root-count bound
all pass across 300 random trials each in both reference implementations.
Expanding $\prod_{a=0}^{16}(x-a)$ directly reproduces $x^{17}-x$
coefficient-by-coefficient, and its $x^1$ coefficient (which is $16$,
i.e.\ $-1 \bmod 17$) matches a direct computation of $16! \bmod 17$,
confirming Wilson's theorem. Exhaustively evaluating 30 random bivariate
polynomials of degree $d \in \{1,\ldots,5\}$ over all $17^2$ points of
$\mathbb{F}_{17}^2$ shows the vanishing frequency never exceeds $d/17$,
matching Schwartz--Zippel.
\end{examplebox}

\section{Further Reading}
\begin{itemize}[leftmargin=1.5em]
  \item AtCoder Algorithm Lectures index (source of this curriculum):
        \url{https://info.atcoder.jp/entry/algorithm_lectures/index}
  \item AtCoder ABC137 F, ``Polynomial Construction'' (an application of
        the $x^p-x$ factorization):
        \url{https://atcoder.jp/contests/abc137/tasks/abc137_f}
  \item Wikipedia, ``Schwartz--Zippel lemma'':
        \url{https://en.wikipedia.org/wiki/Schwartz\%E2\%80\%93Zippel_lemma}
\end{itemize}

\end{document}
